
% --------------------
\section{前言}
    
本文是笔者在讲授《具体数学》部分相关内容时整理的讲义。然而，阅读《具体数学》、完成其部分习题本身是作为课下任务去布置的。在课上，笔者希望能够传达一些书上没有但是在 OI 中较为常见的方法和技巧。

由于时间和难度的原因，讲义无法面面俱到。由于语言的局限，讲义中可能有许多无法解释清楚的地方，也无法完美地传达笔者的所思所想。不过，笔者尽量补齐了课上讲过和希望讲的所有内容，并力图读者\textbf{在 AI 工具的辅助下}能够清楚和深入地理解本讲义的内容与其背后的\textbf{思路}。

本讲义并不完全按照学习的顺序来编写。例如，几乎每道题我们都给出了生成函数的方法，而生成函数本身要到最后才学习。因此，对于这类内容（以 \pending 标记），初学时应当暂时跳过，等到学完相关知识后回头再看。还有一些较为困难的内容（以 * 标记），初学时可能难以完全掌握，可以量力而为。\graffiti{全部跳过。}

《具体数学》是一本非常经典的著作，涵盖了非常丰富的内容。从本书中可以学到很多其他地方难以找到的方法和技巧，也可以从书中看到大师是如何思考问题的。另外，阅读本书可以很好地帮助 OIer 培养离散数学方面的“数学成熟度”。冯·诺依曼曾经说过，“In mathematics you don't understand things. You just get used to them.”（在数学中，你并不理解事物，你只是习惯了它们）。习惯数学是非常重要的！《具体数学》直到 5.4 为止的内容都是对 OI 非常有用的。再往后的内容中，有些与 OI 关联不大，有些会在其他地方学到。但是，读一读仍然会有很大收获。

对于离散数学的初学者来说，《离散数学及其应用》也是一本不错的入门教材。笔者强烈建议至少把这本书的前两章看一看，其中有很多重要且基础的数学语言。后续的内容也都和 OI 比较相关的内容，只不过相对比较基础，只需浏览一下即可。

在讲授本课程期间，笔者让同学们每周上课前自行阅读《具体数学》的对应章节并完成习题。在上课时，仅仅快速地过一遍教材并强调重点，然后开始讲解其他内容。在做书上的习题时，笔者认为不应该查阅前面的正文，以保证对所学知识的记忆。原则上，习题应该全部完成，但为了减轻工作量，只挑选选了一部分有代表性的题目作为书面任务。\graffiti{我上学的时候从不做选做任务。}此外，还有一些补充的 OI 题目。数学习题应当与 OI 题按同样方式对待。笔者期望的课下任务量是每周 8-10 个小时（包含阅读和做题）。

在遇到努力思考但仍然不会的题目时，有以下几种解决办法：
    \begin{itemize}
        \item （对《具体数学》）翻到附录 A 查看答案
        \item （对 OI 题）翻到题解或者网上搜索相关题目的题解
        \item 询问同学/老师/群友
            \begin{itemize}
                \item 理想情况下，你应该表达清楚卡住的点（比如已经尝试过什么、到哪里做不下去了、或者干脆一点头绪没有等等），然后被提问者需要努力提供一个合适的提示。（但是不一定能及时回复）
                \item 而现实中这样沟通的成本往往较高，且问题未必能得到及时的解答。
            \end{itemize}
        \item 询问 AI：AI 工具的出现让上面的问题几乎不复存在了。可以让 AI 来充当上一条中的被提问者。只不过，你需要表达清楚你想让 AI 做的事情，并且需要自行判断 AI 回答的正确性。
    \end{itemize}
现在还可以让 AI 批改自己的作业。参考提示词如下：帮我批改一下这份作业，题目是 xxx（后附作业内容）。另外，规范地使用 Markdown 或 Latex 是非常重要的，你可以向 AI 询问有哪些基本的规范，也可以通过 AI 来学习各种使用技巧。

\graffiti{什么时候智械危机？}